% !TEX root = template.tex
% All numbers are final: single §0.7 test session of 2026-07-22
% (reports/final/) + notebook 06 analysis (paired bootstrap N=10000, seed 42;
% tables in report/tables/, summarized in PROJECT.md Part III).

\section{Results}
\label{sec:results}

\textbf{Statistical protocol.} With $11$ test traces the \emph{trace}, not the
window, is the unit of analysis: the windows of one recording share a channel
and are not independent evidence. Differences are assessed by a trace-level
\emph{paired} bootstrap ($N=10^4$) on \emph{accuracy}; macro-F1 is descriptive
only, since just $\sim$2.8\% of resamples contain all eight classes. Bootstrap
and seed variance are never merged, and since $\sim$10 comparisons share the
same traces, C1--C2 and C1--C3 are primary and the rest exploratory.

\subsection{Main results}
\label{sec:res_main}

\tab{tab:test} reports every configuration with its paired accuracy difference
against our model C1. C1 transfers best (macro-F1 $0.804$, seed-stable to
$0.5$ points) and no variant resolves an improvement over it: SupCon falls
$7.8$ and the GRL $15.8$ points below it, while the largest single effect in
the table is not an objective but the backbone, worth $23.5$ points at equal
recipe.

\begin{table}[t]
\centering
\caption{Main results: validation macro-F1 (checkpoint selection), test
macro-F1 and accuracy on the held-out session, and the paired accuracy
difference vs.\ C1 with its 95\% trace-level bootstrap CI ($\ast$ = interval
excludes zero). Top: primary rotation (leave-lab-out, 8 classes, 11 traces).
Bottom: second rotation (15 traces) and the 5-class SHARP anchor on its own
protocol, neither comparable to the top block.}
\label{tab:test}
\resizebox{\columnwidth}{!}{%
\begin{tabular}{llcccl}
\hline
\textbf{Config} & \textbf{Objective} & \textbf{Val F1} & \textbf{Test F1} & \textbf{Test acc} & \textbf{$\Delta$acc vs C1 [95\% CI]} \\
\hline
C1 & CE (our model) & 0.887 & \textbf{0.804} & \textbf{0.805} & --- \\
C1$'$ (s43) & CE, seed 43 & 0.878 & 0.799 & 0.800 & $-0.005$ \\
C2 & CE + GRL & 0.842 & 0.601 & 0.647 & $-0.158\;[-0.323,+0.034]$ \\
C2$'$ (s43) & CE + GRL, s43 & 0.787 & 0.562 & 0.600 & $-0.205$ \\
C3 & SupCon (probe) & 0.819 & 0.729 & 0.727 & $-0.078\;[-0.128,-0.030]^{\ast}$ \\
C3-ft & SupCon + CE ft & 0.818 & 0.706 & 0.720 & $-0.085\;[-0.221,+0.058]$ \\
C1-sharplike & CE, SHARP net & 0.738 & 0.543 & 0.569 & $-0.235\;[-0.399,-0.092]^{\ast}$ \\
\hline
C1-s6out & CE, living-out & 0.776 & 0.684 & 0.743 & (own rotation) \\
C0 & SHARP anchor (5-cls) & 0.892 & 0.605 & 0.612 & (own split) \\
\hline
\end{tabular}}
\end{table}

\subsection{Our model vs.\ the reference design}
\label{sec:res_baseline}

C0 is an \emph{anchor}, not a benchmark reproduction: it follows the published
architecture and majority-vote fusion on the paper's five-class core, but
holds out $20\%$ of the training capture set for early stopping, which the
reference implementation does not. It therefore supports no claim of the form
``we reproduced $X\%$'', only that our apparatus reproduces the
\emph{behaviour} of the prior approach, degrading to macro-F1 $0.605$ over the
held-out capture sets much as the original work reports. Since C0 differs from
our model in protocol, class set, fusion \emph{and} architecture at once, the
like-for-like comparison is the ablation \emph{C1-sharplike}: the reference
network inside the exact C1 recipe, only the backbone changed. It loses $23.5$
accuracy points (test macro-F1 $0.543$ vs.\ $0.804$), CI $[-0.399,-0.092]$ ---
``deep beats near-linear-wide in this recipe'', on one axis only. Its role is
to close the objection that our CE baseline is a strawman, so that the nulls
below are measured against a strong control.

\subsection{Supervised contrastive learning does not transfer better}
\label{sec:res_supcon}

This is the comparison the study was primarily designed to answer, and the
only one the bootstrap resolves. C3's contrastive phase converged cleanly ---
its pre-committed checkpoint grid spans $0.7$ validation points
($0.819/0.815/0.812$ at epochs $40/50/60$), so the selected checkpoint sits on
a plateau --- and on test it lands $7.8$ accuracy points below C1, CI
$[-0.128,-0.030]$.

Since a frozen linear probe could in principle understate a non-linear
representation, we attacked that explanation from four directions.
\emph{Readouts:} on the same frozen validation features C3 scores $0.819$
(linear), $0.805$ (kNN) and $0.718$ (nearest-class mean) in macro-F1 against
C1's $0.884$, $0.856$ and $0.889$; C1 is robust to the readout while C3 loses
under all three, and although kNN beats the class-mean readout on C3 --- so
non-linear structure is genuinely present --- it still does not reach C3's own
linear probe. \emph{Fine-tuning:} unfreezing everything and training with CE
from the SupCon initialization (C3-ft) was pre-registered with the hypothesis
``comparable to C1 within one point'' and falsified it, settling at $0.818$ on
validation --- its own linear-probe ceiling --- and $0.720$ on test.
\emph{Complementarity:} concatenated C1$\oplus$C3 features score $0.868$ on
validation, \emph{below} the same-loss ensemble control C1$\oplus$C1$'$
($0.888$): SupCon is a worse partner than a second CE encoder.
\emph{Geometry:} the left column of \fig{fig:tsne} shows why a linear readout
struggles --- C3's low-motion classes (L, S and the empty room) chain into one
filament instead of three separable blobs --- and the same diagnostic on C3-ft
shows the fine-tune \emph{un-chaining} that geometry, turning the SupCon
encoder back into C1 rather than building on it. Counted honestly this is two families of evidence,
not four instruments: the readouts share cached features and agree partly by
construction, while the fine-tune is a different path and carries the
confirmation. Both point the same way --- the SupCon features are less
informative about the activity, and no head, linear or not, recovers what they
do not encode.

\subsection{The adversary costs accuracy and stability}
\label{sec:res_grl}

C2 loses $15.8$ accuracy points on test. The direction is stable --- it holds
on both seeds (seed-43 gap $20.0$ points) and on validation --- but the
interval $[-0.323,+0.034]$ crosses zero: on $11$ traces even a $16$-point gap
is unresolved. The replicates add a second observation: the
same-configuration spread is
$0.5$ accuracy points for C1 against $4.7$ for C2, so the adversary appears to
cost stability as well as accuracy. A linear probe on C2's \emph{frozen}
features reaches $0.720$ test accuracy, $7$ points above C2's own trained head
($0.647$), against a probe--head gap of $+0.002$ for C1: much of the damage
sits in the classifier head, not in the representation.

\subsection{Why: the adversary has no target}
\label{sec:res_diag}

We answer the natural question --- why an invariance objective fails on a
transfer benchmark --- by probing the \emph{frozen} encoders along the dashed
path of \fig{fig:learning}. On a trace-disjoint inner split of the training
traces we fit linear probes for the activity (positive control) and for every
domain attribute, each against its own majority baseline. The control
saturates ($1.000$ against a $0.197$ baseline) while \emph{every} domain
attribute sits at or below its own: capture set $-0.090$, direct path
$-0.098$, monitor position $-0.086$, and the environment and subject probes
collapse to the constant majority predictor. None of them is linearly readable
from the features, on precisely the data an in-training adversary reads. The
same probes on the C2 and C3 encoders, on all three contrastive checkpoints
and on the second rotation --- five encoders in total --- return the same
verdict, the largest positive deviation anywhere being $+0.015$; the right
column of \fig{fig:tsne} is the geometric counterpart, capture-set colours
uniformly mixed inside every activity cluster.

The implication is structural. A GRL removes domain information its
discriminator can read; here there is nothing to read, so it can only inject
gradient noise --- which is what the training monitors showed (on both seeds
the domain head never left the majority-class floor) and what the test rows
confirm. The root cause is the environment incidence of \secref{sec:model}:
with the bedroom in five capture sets and one set for each other room, any
rotation trains either on a single environment or on one environment plus a
single session, so environment invariance is barely definable on this training
distribution --- which is also why C4 was closed without training, as
evidence-based triage rather than a claim that the GRL is useless in general.
The probe alone shows only that the domain is not \emph{linearly} decodable
from traces the encoder has seen; the structural argument is what generalizes
the verdict.

\begin{figure}[t]
\centering
\includegraphics[width=0.88\columnwidth]{embeddings_c1_vs_c3.pdf}
\caption{t-SNE~\cite{vanderMaaten-2008} of train features; top row C1 (CE),
bottom row C3 (SupCon). \emph{Left, by activity:} both form tight clusters,
but C3's L/S/E clusters chain into one filament instead of three separable
blobs --- the geometric counterpart of its lower linear-probe score.
\emph{Right, the same points by capture set:} colours are uniformly mixed
inside every cluster, in both encoders --- no domain sub-structure for an
adversary to remove.}
\label{fig:tsne}
\end{figure}

\subsection{Error structure and limitations}
\label{sec:res_err}

Checkpoint selection was blind to C, E and S, and we had pre-declared the risk
that those classes would underperform; the opposite holds, the val-blind
classes being the \emph{best} ($0.904$ mean F1 against $0.744$), so class
difficulty in the target environment, not selection blindness, drives
per-class error. The errors themselves are
physically sensible (\fig{fig:confusion}): C1 confuses mainly
walking$\leftrightarrow$running and H$\to$J, pairs with adjacent Doppler
signatures, and recognizes the empty room perfectly --- confusions between
physically similar activities, i.e.\ genuine motion features rather than room
artifacts. Calibration
dissociates from accuracy: C1 is systematically \emph{under}confident (ECE
$0.243$, an expected effect of label smoothing) while the GRL model is the
best calibrated ($0.070$) and the least accurate, so any deployment consuming
confidence scores needs temperature scaling.

\textbf{Scope and limitations.} The same session also evaluated two further
pre-registered arms that we do not discuss for lack of space --- test-time
adaptation (AdaBN~\cite{Li-2017}, T3A~\cite{Iwasawa-2021}) and amplitude
augmentation --- neither of which resolved an improvement over C1 (best
$+0.007$, with an interval crossing zero). Beyond those, model selection is
in-domain ($15\%$ of the training traces), so the selected checkpoints may be
sub-optimal out of domain; each held-out set is a single session, so
``environment'' is confounded with subject, monitor position and day; and at
$11$ and $15$ traces a percentile bootstrap separates two or three performance
tiers, not a dozen individually ranked rows.

\begin{figure}[t]
\centering
\includegraphics[width=0.44\columnwidth]{confusion_c1_test.pdf}
\caption{Confusion matrix of C1 on the held-out laboratory session
(row-normalized).}
\label{fig:confusion}
\end{figure}
